% Make sure that any definitions we do in this file stay local.
\begingroup

\newlength\msgSelfWidth
\newlength\msgSelfHeight
\newlength\msgSelfOffset
\newlength\activationWidth
\newlength\firstStateOffset

\msgSelfWidth=0.9cm\relax
\msgSelfHeight=0.25cm\relax
\msgSelfOffset=0.15cm\relax
\activationWidth=0.14cm\relax
\firstStateOffset=0.08cm\relax

\newsavebox\SplitDrawingLocal
\newsavebox\SplitDrawingRemote

\newcommand\saveframedtikz[3][]{%
  \sbox{#2}{%
    \fbox{%
      \hspace{1mm}%
      \begin{tikzpicture}[#1]#3\end{tikzpicture}%
      \hspace{1mm}%
    }%
  }%
}

\tikzset{
  >=Stealth,
  head/.style={draw, minimum width=1.15cm,
    minimum height=0.42cm, inner sep=1pt, font=\scriptsize},
  state/.style={draw, fill=blue!5, align=center, inner sep=2pt,
    font=\scriptsize, text width=1.45cm, xshift = 2pt, anchor = west},
  flag/.style={draw=orange!65!black, fill=orange!14},
  inactive/.style={state, fill=black!12},
  lifeline/.style={dashed, black!55},
  activation/.style={draw, fill=black!8},
  msg/.style={->},
  create/.style={->, dashed},
  return/.style={->, dashed},
  lab/.style={font=\scriptsize},
  lab above/.style={lab,above,anchor=base,yshift=2.4pt},
  lab right/.style={lab,right,anchor=west,yshift=-.8pt},
}

\saveframedtikz\SplitDrawingLocal{
  \node[head] (p) at (0,0) {$P$};
  %\draw[lifeline] (p.south) -- ++(0,-5.3);

  % Define three points of the activation rectangle: upper right, lower
  % right and lower left.
  \path (p.south) ++(0, 0.2pt)
    ++( 0.5\activationWidth,  0) coordinate (act ur)
    ++( 0.0,                 -5) coordinate (act lr)
    ++(-1.0\activationWidth,  0) coordinate (act ll);

  % Draw the activation rectangle.
  \draw[activation] (act ur) rectangle (act ll);

  % Place nodes along the activation rectangle.
  \node[state] (c0) at ($(act ur)!.07!(act lr)$) {$c:S_1;S_2$};
  \node[state] (c1) at ($(act ur)!.35!(act lr)$) {$c_1:S_1$};
  \node[inactive, anchor = north west, yshift=.4pt] (c2 inactive)
    at ($(act ur)!(c1.south)!(act lr)$) {$c_2:S_2$\\inactive};
  \node[state] (c2) at ($(act ur)!.75!(act lr)$) {$c_2:S_2$};

  % Project (c1.north) onto the activation rectangle and draws the arrow
  % with the dimensions in \msgSelfWidth, \msgSelfHeight, \msgSelfOffset.
  \draw[msg] ($(act ur)!(c1.north)!(act lr)$)
    ++(0, \msgSelfOffset) ++(0, \msgSelfHeight)
    -- node[lab above] {$\CodeOp{lsplit}\,c$} ++(\msgSelfWidth,0)
    -- ++(0, -\msgSelfHeight)
    -- ++(-\msgSelfWidth, 0);

  % Repeat for the "use c1" arrow.
  \draw[msg] ($(act ur)!(c2.north)!(act lr)$)
    ++(0, \msgSelfOffset) ++(0, \msgSelfHeight)
    -- node[lab above] {use $c_1$} ++(\msgSelfWidth,0)
    -- ++(0, -\msgSelfHeight)
    -- ++(-\msgSelfWidth, 0);

  % Repeat for the "use c2" arrow. The arrow point is set 95% along the
  % activation rectangle.
  \draw[msg] ($(act ur)!.95!(act lr)$)
    ++(0, \msgSelfHeight)
    -- node[lab above] {use $c_2$} ++(\msgSelfWidth,0)
    -- ++(0, -\msgSelfHeight)
    -- ++(-\msgSelfWidth, 0);
}

\saveframedtikz\SplitDrawingRemote{
  \node[head] (p) {$P$};
  \node[head, flag, below right=0.37cm and 2.5cm of p] (z) {$\varphi\,z$};
  \node[head,       below right=0.61cm and 2.5cm of z] (q) {$Q$};

  % Coordinates for the upper edge of p's first activation.
  \path (p.south) ++(0, .2pt)
    + (-.5\activationWidth,  0) coordinate (p1 ul)
    ++( .5\activationWidth,  0) coordinate (p1 ur)
    ++(                  0, -1) coordinate (p er);

  % The initial c channel.
  \path (p1 ur) ++(0, -\firstStateOffset)
    node[state, anchor=north west] (c) {$c:S_1;S_2$};

  % p's lifeline
  \draw[lifeline] (p.south) -- ++(0, -5) coordinate (p end);

  % Define three points of p's first activation rectangle: upper right,
  % lower right and lower left.
  \path (p1 ur)
    ++( 0.0,                 -2) coordinate (p1 lr)
    ++(-1.0\activationWidth,  0) coordinate (p1 ll);
  \draw[activation] (p1 ur) rectangle (p1 ll);

  % q's activation rectangle
  \path ($(q.north)!(p end)!(q.south)$)
    ++(0, .5)
    +( .5\activationWidth, 0) coordinate (q lr)
    +(-.5\activationWidth, 0) coordinate (q ll);
  \path (q.south) +(.5\activationWidth, .2pt) coordinate (q ur);
  \draw[activation] (q ur) rectangle (q ll);

  % Define coordinates that define the right edge of z's activation
  % rectangle.
  \path (z.south) ++(0, .2pt)
    ++(0.5\activationWidth, 0) coordinate (z ur)
    ++(0, -1) coordinate (z er);

  % actions and states in Q
  \path (q ur) ++(0, -\firstStateOffset)
    node[state, anchor=north west] (c1 at q) {$c_1:S_1;\CodeTy{Drop}$};
  \draw[msg] ($(q ur)!(c1 at q.south)!(q lr)$)
    ++ (0, -\msgSelfOffset)
    -- ++(\msgSelfWidth,0)
    -- node[lab right] {use $c_1$} ++(0, -\msgSelfHeight)
    -- ++(-\msgSelfWidth, 0) coordinate (q use end);
  \path (q use end) ++(0, -\msgSelfOffset)
    node[state, anchor=north west] (c1 drop) {$c_1:\CodeTy{Drop}$};

  % the "drop" arrow
  \path ($(q ur)!(c1 drop.south)!(q lr)$)
    ++(-\activationWidth,-\msgSelfOffset)
    coordinate (drop origin);
  \draw[msg] (drop origin) -- node[lab above] {$\CodeOp{drop}\,c_1$}
    ($(z ur)!(drop origin)!(z er)$) coordinate (drop target);

  % States for the synchronization flag.
  \path (z ur) ++(0, -\firstStateOffset)
    node[state, flag, anchor = north west] {$z\mapsto\texttt{drop}$};
  \path (drop target) ++(0, -\msgSelfOffset)
    node[state, flag, anchor = north west] (z acq) {$z\mapsto\texttt{acq}$};

  % Define the remaining coordinates for z's activation rectangle.
  \path (z ur) ++(-\activationWidth, 0) coordinate (z ul);
  \path ($(z ur)!(z acq.south)!(z er)$)
    ++(0, -\msgSelfOffset) coordinate (z lr)
    ++(-\activationWidth, 0) coordinate (z ll);
  \draw[activation, flag] (z ur) rectangle (z ll);

  % Define three points of p's second activation rectangle: lower right,
  % lower left and upper right.
  \path (p end) +( .5\activationWidth, 0) coordinate (p2 lr)
                +(-.5\activationWidth, 0) coordinate (p2 ll);
  \coordinate (p2 ur) at ($(p2 lr)!(z ll)!(p1 lr)$);
  \draw[activation] (p2 ur) rectangle (p2 ll);

  % z's creation arrow
  \draw[create]
    ($(p1 ur)!(z.west)!(p1 lr)$) coordinate (rsplit origin)
    -- node[lab above, near end] {$\CodeOp{rsplit}\,c$}
       node[lab, below, yshift=2.2pt, near end] {create $z$} (z.west);

  % c1 and c2 immediately after the rsplit
  \path (rsplit origin) ++(0, -\msgSelfOffset)
    node[state, anchor = north west] (c1 at p) {$c_1:S_1;\CodeTy{Drop}$};
  \node[state, anchor = north west, yshift=.4pt] (c1 at p)
    at ($(p1 ur)!(c1 at p.south)!(p er)$) {$c_2:\mathrlap{\CodeTy{Acq};S_2}\hphantom{S_1;\CodeTy{Drop}}$};

  % Remember the coordinate where the fork arrow starts.
  \coordinate (fork q) at ($(p1 ur)!(q.west)!(p1 lr)$);

  % q's fork arrow
  \draw[create] (fork q) -- (q);
  \path (fork q) -- node[lab above, near end] {$\CodeOp{fork}$}
                    node[state, below=2pt, near end] {$c_1 : S_1;\CodeTy{Drop}$} (q);

  % the "acquire" arrow
  \draw[msg] (p1 lr) --
    node[lab above] {$\CodeOp{acquire}\,c_2$}
    ($(z ul)!(p1 lr)!(z ll)$);
  % the "return" arrow
  \draw[return] (z ll) -- node[lab above] {return} (p2 ur);

  \path (p2 ur) ++(0, -\msgSelfOffset)
    node[state, anchor=north west] (c2 ret) {$c_2:S_2$};
  \draw[msg] ($(p2 ur)!(c2 ret.south)!(p2 lr)$) ++(0,-\msgSelfOffset)
    -- ++(\msgSelfWidth,0)
    -- node[lab right] {use $c_2$} ++(0,-\msgSelfHeight)
    -- ++(-\msgSelfWidth,0);
}

\begin{figure}[t]
  \begin{subfigure}[t]{\wd\SplitDrawingLocal}
    \usebox\SplitDrawingLocal
    \caption{Local borrow}
    \label{fig:borrow-lifecycle-msc-local}
  \end{subfigure}\hfill
  \begin{subfigure}[t]{\wd\SplitDrawingRemote}
    \usebox\SplitDrawingRemote
    \caption{Remote borrow}
    \label{fig:borrow-lifecycle-msc-remote}
  \end{subfigure}
  \caption{Message-sequence view of local and remote borrows. Channel state is
    shown as annotations next to the relevant process or flag lifelines. A local
    split keeps the channel state with the original process lifeline; a remote
    split creates a flag lifeline $\varphi\,z$ and forks a process for the
    borrowed prefix.}
  \label{fig:borrow-lifecycle-msc}
\end{figure}

% Closes the group opened at the start of the file.
\endgroup

%%% Local Variables:
%%% mode: LaTeX
%%% TeX-master: "popl2027"
%%% End:
